%% Modelo LaTeX para dissertações e teses do PPGT-UnB
%% Baseado no template oficial do Programa de Pós-Graduação em Transportes
%% Universidade de Brasília - Faculdade de Tecnologia
%%
%% Este arquivo serve como exemplo completo de uso da classe UnB-PPGT
%% Para mais informações, consulte o README.md

%% Opções da classe:
%% - mestrado (padrão) | doutorado | qualificacao
%% - onehalfspacing (padrão) | singlespacing | doublespacing  
%% - 12pt (padrão) | 10pt | 11pt
%% - hyperref (padrão) | impressao

\documentclass[mestrado,onehalfspacing,12pt,hyperref]{UnB-PPGT}

%% ========================================================================
%% METADADOS DO DOCUMENTO
%% ========================================================================

%% Título e subtítulo
\titulo{Análise da Mobilidade Urbana Sustentável em Cidades de Médio Porte}
\subtitulo{Um Estudo de Caso da Região Metropolitana de Brasília}

%% Autor
\autor{João}{Silva Santos}

%% Orientação
\orientador{Prof. Dr. Maria Fernanda Oliveira}{Universidade de Brasília}
\coorientador{Prof. Dr. Carlos Eduardo Lima}{Universidade Federal de Goiás}

%% Data da defesa
\diamesano{15}{3}{2024}

%% Coordenador do programa
\coordenador{Prof. Dr. Ana Paula Costa}{Universidade de Brasília}

%% Banca examinadora
\membrobancafunc{Prof. Dr. Maria Fernanda Oliveira}{Universidade de Brasília}{Orientadora}
\membrobancafunc{Prof. Dr. Carlos Eduardo Lima}{Universidade Federal de Goiás}{Coorientador}
\membrobancafunc{Prof. Dr. Roberto Silva}{Universidade de Brasília}{Examinador Interno}
\membrobancafunc{Prof. Dr. Ana Beatriz Santos}{Universidade Federal de Minas Gerais}{Examinadora Externa}
\membrobancafunc{Prof. Dr. Pedro Henrique Costa}{Instituto de Pesquisa Econômica Aplicada}{Examinador Externo}

%% Palavras-chave (para resumo e ficha catalográfica)
\palavraschave[Mobilidade urbana. Sustentabilidade. Transporte público. Planejamento urbano. Brasília.]{mobilidade urbana, sustentabilidade, transporte público, planejamento urbano, Brasília}

%% Keywords (para abstract e ficha catalográfica)
\keywords[Urban mobility. Sustainability. Public transportation. Urban planning. Brasília.]{urban mobility, sustainability, public transportation, urban planning, Brasília}

%% ========================================================================
%% PACOTES ADICIONAIS (se necessário)
%% ========================================================================

%% Pacotes para matemática avançada
\usepackage{amsmath,amssymb,amsthm}

%% Pacotes para tabelas e figuras
\usepackage{booktabs}
\usepackage{multirow}
\usepackage{subcaption}

%% Pacotes para algoritmos (se necessário)
% \usepackage{algorithm}
% \usepackage{algorithmic}

%% Pacotes para código fonte (se necessário)
% \usepackage{listings}
% \usepackage{xcolor}

%% ========================================================================
%% CONFIGURAÇÕES ADICIONAIS
%% ========================================================================

%% Definição de novos comandos (exemplos)
\newcommand{\sigla}[1]{\textsc{#1}}
\newcommand{\software}[1]{\textit{#1}}
\newcommand{\variavel}[1]{\textit{#1}}

%% Configurações para teoremas (se necessário)
\newtheorem{teorema}{Teorema}[chapter]
\newtheorem{definicao}{Definição}[chapter]
\newtheorem{proposicao}{Proposição}[chapter]

%% ========================================================================
%% INÍCIO DO DOCUMENTO
%% ========================================================================

\begin{document}

%% ========================================================================
%% SEÇÃO PRÉ-TEXTUAL
%% ========================================================================

\pretextual

%% Capa (gerada automaticamente)
\capa

%% Folha de rosto (gerada automaticamente)
\folhaderosto

%% Folha de aprovação (gerada automaticamente)
\folhadeaprovacao

%% Ficha catalográfica (gerada automaticamente)
\fichacatalografica

%% Dedicatória (opcional)
\begin{dedicatoria}
Dedico este trabalho aos meus pais, Maria e José, que sempre acreditaram na importância da educação e me apoiaram incondicionalmente durante toda esta jornada acadêmica.

Aos profissionais do transporte público que, diariamente, trabalham para melhorar a mobilidade urbana em nossas cidades.
\end{dedicatoria}

%% Agradecimentos (opcional)
\begin{agradecimentos}
Agradeço primeiramente a Deus, por me conceder força e sabedoria durante todo o percurso desta pesquisa.

À minha orientadora, Prof. Dr. Maria Fernanda Oliveira, pela paciência, dedicação e valiosas contribuições que tornaram possível a realização deste trabalho. Sua experiência e conhecimento foram fundamentais para o desenvolvimento desta dissertação.

Ao meu coorientador, Prof. Dr. Carlos Eduardo Lima, pelas importantes sugestões e pelo apoio técnico durante as análises estatísticas.

Aos professores do Programa de Pós-Graduação em Transportes da UnB, pelos ensinamentos transmitidos e pela formação acadêmica de excelência.

Aos colegas do PPGT, especialmente aos membros do grupo de pesquisa em Mobilidade Urbana, pelas discussões enriquecedoras e pelo companheirismo.

À Coordenação de Aperfeiçoamento de Pessoal de Nível Superior (CAPES), pela concessão da bolsa de estudos que viabilizou a dedicação integral a esta pesquisa.

Aos funcionários das empresas de transporte público e órgãos governamentais que colaboraram fornecendo dados essenciais para esta pesquisa.

À minha família e amigos, pelo apoio emocional e compreensão durante os momentos de ausência.

A todos que, direta ou indiretamente, contribuíram para a realização deste trabalho.
\end{agradecimentos}

%% Resumo
\begin{resumo}
A mobilidade urbana sustentável representa um dos principais desafios das cidades contemporâneas, especialmente em regiões metropolitanas de países em desenvolvimento. Este trabalho apresenta uma análise abrangente da mobilidade urbana na Região Metropolitana de Brasília, com foco na sustentabilidade dos sistemas de transporte. O objetivo principal é avaliar as condições atuais de mobilidade e propor estratégias para o desenvolvimento de um sistema de transporte mais sustentável e eficiente.

A metodologia adotada combina análise quantitativa de dados de mobilidade, pesquisas de campo com usuários do transporte público e análise qualitativa de políticas públicas. Foram coletados dados de deslocamentos, tempos de viagem, custos operacionais e impactos ambientais dos diferentes modais de transporte. A pesquisa abrangeu o período de 2020 a 2023, permitindo uma análise temporal das transformações no padrão de mobilidade.

Os resultados indicam que o transporte individual motorizado ainda representa 65\% dos deslocamentos na região, enquanto o transporte público atende apenas 25\% da demanda. O transporte não motorizado (caminhada e bicicleta) corresponde a 10\% dos deslocamentos. A análise revelou deficiências significativas na integração entre os diferentes modais, baixa frequência do transporte público em horários de pico e infraestrutura inadequada para pedestres e ciclistas.

O estudo propõe um conjunto de estratégias integradas para promover a mobilidade sustentável, incluindo: (i) expansão e modernização do sistema de transporte público; (ii) implementação de corredores exclusivos para ônibus; (iii) desenvolvimento de infraestrutura cicloviária; (iv) políticas de desestímulo ao uso do automóvel; e (v) integração tarifária e operacional entre os modais.

As contribuições desta pesquisa incluem uma metodologia replicável para avaliação da mobilidade urbana, um diagnóstico detalhado da situação atual da Região Metropolitana de Brasília e um conjunto de recomendações práticas para gestores públicos. Os resultados demonstram que a transição para um sistema de mobilidade mais sustentável requer investimentos coordenados em infraestrutura, tecnologia e políticas públicas integradas.
\end{resumo}

%% Abstract
\begin{abstract}
Sustainable urban mobility represents one of the main challenges of contemporary cities, especially in metropolitan regions of developing countries. This work presents a comprehensive analysis of urban mobility in the Metropolitan Region of Brasília, focusing on the sustainability of transportation systems. The main objective is to evaluate current mobility conditions and propose strategies for developing a more sustainable and efficient transportation system.

The adopted methodology combines quantitative analysis of mobility data, field surveys with public transport users, and qualitative analysis of public policies. Data on trips, travel times, operational costs, and environmental impacts of different transport modes were collected. The research covered the period from 2020 to 2023, allowing for a temporal analysis of transformations in mobility patterns.

Results indicate that individual motorized transport still represents 65\% of trips in the region, while public transport serves only 25\% of demand. Non-motorized transport (walking and cycling) accounts for 10\% of trips. The analysis revealed significant deficiencies in integration between different modes, low frequency of public transport during peak hours, and inadequate infrastructure for pedestrians and cyclists.

The study proposes a set of integrated strategies to promote sustainable mobility, including: (i) expansion and modernization of the public transport system; (ii) implementation of exclusive bus corridors; (iii) development of cycling infrastructure; (iv) policies to discourage automobile use; and (v) fare and operational integration between modes.

The contributions of this research include a replicable methodology for urban mobility assessment, a detailed diagnosis of the current situation in the Metropolitan Region of Brasília, and a set of practical recommendations for public managers. The results demonstrate that the transition to a more sustainable mobility system requires coordinated investments in infrastructure, technology, and integrated public policies.
\end{abstract}

%% Lista de figuras (gerada automaticamente)
\listoffigures

%% Lista de tabelas (gerada automaticamente)
\listoftables

%% Lista de siglas e abreviações
\begin{siglas}
\item[ANTP] Associação Nacional de Transportes Públicos
\item[BRT] \textit{Bus Rapid Transit}
\item[CAPES] Coordenação de Aperfeiçoamento de Pessoal de Nível Superior
\item[DF] Distrito Federal
\item[DETRAN] Departamento de Trânsito
\item[IBGE] Instituto Brasileiro de Geografia e Estatística
\item[IPEA] Instituto de Pesquisa Econômica Aplicada
\item[PDTU] Plano Diretor de Transporte Urbano
\item[PPGT] Programa de Pós-Graduação em Transportes
\item[RIDE] Região Integrada de Desenvolvimento do Entorno
\item[SEMOB] Secretaria de Mobilidade
\item[SIG] Sistema de Informações Geográficas
\item[STPC] Sistema de Transporte Público Coletivo
\item[UnB] Universidade de Brasília
\item[VLT] Veículo Leve sobre Trilhos
\end{siglas}

%% ========================================================================
%% SEÇÃO TEXTUAL
%% ========================================================================

\textual

%% Capítulos do trabalho
%% ========================================================================
%% CAPÍTULO 1 - INTRODUÇÃO
%% ========================================================================

\capitulo{intro}{Introdução}

A mobilidade urbana constitui um dos principais desafios das cidades contemporâneas, especialmente em regiões metropolitanas de países em desenvolvimento. O crescimento populacional acelerado, associado à expansão urbana desordenada, tem gerado problemas complexos relacionados ao transporte de pessoas e mercadorias \cite{vasconcellos2012}.

No Brasil, as questões de mobilidade urbana ganharam destaque nas últimas décadas devido ao aumento significativo da frota de veículos particulares e à inadequação dos sistemas de transporte público. Segundo dados do \sigla{IBGE} \cite{ibge2022}, a frota nacional de veículos cresceu 300\% entre 1990 e 2020, enquanto a população urbana aumentou apenas 60\% no mesmo período.

\section{Contextualização}

A Região Metropolitana de Brasília, composta pelo Distrito Federal e municípios do entorno, apresenta características peculiares que influenciam diretamente os padrões de mobilidade urbana. O planejamento original da capital, concebido por Lucio Costa na década de 1950, priorizou o transporte individual motorizado, resultando em uma estrutura urbana dispersa e dependente do automóvel \cite{costa2018}.

Atualmente, a região enfrenta desafios significativos relacionados à mobilidade, incluindo:

\begin{itemize}
    \item Congestionamentos frequentes nos principais corredores viários;
    \item Baixa qualidade do transporte público;
    \item Infraestrutura inadequada para pedestres e ciclistas;
    \item Altos índices de poluição atmosférica;
    \item Segregação socioespacial relacionada à acessibilidade.
\end{itemize}

\section{Justificativa}

A necessidade de estudos sobre mobilidade urbana sustentável na Região Metropolitana de Brasília justifica-se por diversos fatores. Primeiro, a região apresenta um dos maiores índices de motorização do país, com aproximadamente 0,7 veículos por habitante \cite{detran2023}. Segundo, os custos sociais e ambientais do modelo atual de mobilidade são crescentes, incluindo tempo perdido em congestionamentos, acidentes de trânsito e emissões de gases poluentes.

Além disso, a implementação de políticas públicas de mobilidade urbana requer fundamentação técnica e científica adequada. A Lei Federal nº 12.587/2012, que institui a Política Nacional de Mobilidade Urbana, estabelece diretrizes para o desenvolvimento de sistemas de transporte mais sustentáveis e equitativos \cite{brasil2012}.

\section{Problema de Pesquisa}

O problema central desta pesquisa pode ser formulado da seguinte forma: \textbf{Como promover a transição para um sistema de mobilidade urbana mais sustentável na Região Metropolitana de Brasília, considerando as características específicas da região e as demandas da população?}

Esta questão principal desdobra-se em problemas específicos:

\begin{enumerate}
    \item Quais são as características atuais dos padrões de mobilidade na região?
    \item Quais os principais obstáculos para a adoção de modos de transporte mais sustentáveis?
    \item Que estratégias podem ser implementadas para melhorar a sustentabilidade do sistema de transporte?
    \item Como integrar os diferentes modais de transporte de forma eficiente?
\end{enumerate}

\section{Objetivos}

\subsection{Objetivo Geral}

Analisar as condições atuais de mobilidade urbana na Região Metropolitana de Brasília e propor estratégias integradas para o desenvolvimento de um sistema de transporte mais sustentável e eficiente.

\subsection{Objetivos Específicos}

\begin{enumerate}
    \item Caracterizar os padrões atuais de mobilidade urbana na região, identificando os principais modais utilizados e suas características operacionais;
    
    \item Avaliar a qualidade dos serviços de transporte público, considerando aspectos como acessibilidade, frequência, conforto e integração;
    
    \item Analisar os impactos ambientais, sociais e econômicos dos diferentes modais de transporte;
    
    \item Identificar as principais barreiras para a adoção de modos de transporte mais sustentáveis;
    
    \item Propor um conjunto de estratégias integradas para promover a mobilidade urbana sustentável;
    
    \item Desenvolver indicadores para monitoramento e avaliação das políticas de mobilidade urbana.
\end{enumerate}

\section{Hipóteses}

As hipóteses que orientam esta pesquisa são:

\begin{enumerate}
    \item A predominância do transporte individual motorizado na Região Metropolitana de Brasília está relacionada às deficiências do sistema de transporte público e à inadequação da infraestrutura para modos não motorizados;
    
    \item A implementação de estratégias integradas de mobilidade urbana pode reduzir significativamente a dependência do automóvel e melhorar a sustentabilidade do sistema de transporte;
    
    \item A integração física, tarifária e operacional entre os diferentes modais de transporte é fundamental para o sucesso de políticas de mobilidade sustentável.
\end{enumerate}

\section{Estrutura do Trabalho}

Este trabalho está estruturado em quatro capítulos principais:

O \refCap{intro} apresenta a contextualização do problema, justificativa, objetivos e hipóteses da pesquisa.

O \refCap{metodologia} descreve os procedimentos metodológicos adotados, incluindo a caracterização da área de estudo, métodos de coleta e análise de dados.

O \refCap{resultados} apresenta e discute os resultados obtidos, incluindo a caracterização dos padrões de mobilidade, avaliação dos sistemas de transporte e proposição de estratégias.

O \refCap{conclusao} sintetiza as principais conclusões, contribuições da pesquisa e recomendações para trabalhos futuros.

\section{Delimitação do Estudo}

Esta pesquisa tem como recorte espacial a Região Metropolitana de Brasília, incluindo o Distrito Federal e os municípios do entorno que fazem parte da \sigla{RIDE}. O recorte temporal abrange o período de 2020 a 2023, permitindo uma análise das transformações recentes nos padrões de mobilidade.

O estudo foca nos deslocamentos de pessoas, não abordando especificamente o transporte de cargas. Embora reconheça-se a importância da logística urbana para a mobilidade sustentável, esta temática não será aprofundada devido às limitações de escopo da pesquisa.
%% ========================================================================
%% CAPÍTULO 2 - METODOLOGIA
%% ========================================================================

\capitulo{metodologia}{Metodologia}

Este capítulo apresenta os procedimentos metodológicos adotados para o desenvolvimento da pesquisa. A metodologia combina abordagens quantitativas e qualitativas, permitindo uma análise abrangente da mobilidade urbana na Região Metropolitana de Brasília.

\section{Caracterização da Pesquisa}

Esta pesquisa caracteriza-se como um estudo de caso exploratório e descritivo, com abordagem mista (quantitativa e qualitativa). Segundo \cite{yin2015}, o estudo de caso é apropriado quando se busca compreender fenômenos contemporâneos complexos em seu contexto real.

A escolha da abordagem mista justifica-se pela necessidade de:
\begin{itemize}
    \item Quantificar padrões de mobilidade através de dados estatísticos;
    \item Compreender percepções e comportamentos dos usuários;
    \item Analisar políticas públicas e sua implementação;
    \item Triangular diferentes fontes de evidência.
\end{itemize}

\section{Área de Estudo}

A área de estudo compreende a Região Metropolitana de Brasília, conforme delimitação da \sigla{RIDE}, estabelecida pela Lei Complementar nº 94/1998. A região inclui:

\begin{itemize}
    \item \textbf{Distrito Federal}: Brasília e demais Regiões Administrativas;
    \item \textbf{Goiás}: Águas Lindas de Goiás, Alexânia, Cidade Ocidental, Cocalzinho de Goiás, Corumbá de Goiás, Cristalina, Formosa, Luziânia, Novo Gama, Padre Bernardo, Pirenópolis, Planaltina, Santo Antônio do Descoberto, Valparaíso de Goiás e Vila Boa;
    \item \textbf{Minas Gerais}: Buritis e Unaí.
\end{itemize}

A \refFig{fig:area-estudo} apresenta a delimitação da área de estudo e suas principais características geográficas.

\begin{figure}[htb]
    \centering
    \includegraphics[width=0.8\textwidth]{img/area-estudo.png}
    \caption{Delimitação da área de estudo - Região Metropolitana de Brasília}
    \label{fig:area-estudo}
\end{figure}

\section{Procedimentos Metodológicos}

A metodologia foi estruturada em cinco etapas principais, conforme apresentado na \refFig{fig:fluxograma-metodologia}:

\begin{figure}[htb]
    \centering
    \includegraphics[width=0.9\textwidth]{img/fluxograma-metodologia.png}
    \caption{Fluxograma da metodologia de pesquisa}
    \label{fig:fluxograma-metodologia}
\end{figure}

\subsection{Etapa 1: Revisão Bibliográfica}

A revisão bibliográfica foi conduzida de forma sistemática, abrangendo:

\begin{itemize}
    \item Conceitos fundamentais de mobilidade urbana sustentável;
    \item Experiências internacionais em sistemas de transporte sustentável;
    \item Políticas públicas de mobilidade urbana no Brasil;
    \item Estudos específicos sobre a Região Metropolitana de Brasília.
\end{itemize}

As bases de dados consultadas incluíram \software{Web of Science}, \software{Scopus}, \software{SciELO} e repositórios institucionais. Os critérios de seleção priorizaram publicações dos últimos 10 anos, com foco em estudos empíricos e revisões sistemáticas.

\subsection{Etapa 2: Coleta de Dados Secundários}

Os dados secundários foram obtidos junto a diferentes instituições:

\begin{itemize}
    \item \textbf{\sigla{IBGE}}: Dados demográficos e socioeconômicos;
    \item \textbf{\sigla{DETRAN}-DF}: Dados da frota de veículos;
    \item \textbf{\sigla{SEMOB}-DF}: Dados do transporte público;
    \item \textbf{\sigla{ANTP}}: Indicadores de mobilidade urbana;
    \item \textbf{Empresas operadoras}: Dados operacionais do transporte público.
\end{itemize}

A \refTab{tab:dados-secundarios} apresenta um resumo dos dados coletados e suas respectivas fontes.

\begin{table}[htb]
    \centering
    \caption{Dados secundários coletados por fonte}
    \label{tab:dados-secundarios}
    \begin{tabular}{lll}
        \toprule
        \textbf{Variável} & \textbf{Fonte} & \textbf{Período} \\
        \midrule
        População & IBGE & 2010-2022 \\
        Frota de veículos & DETRAN-DF & 2020-2023 \\
        Passageiros transportados & SEMOB-DF & 2020-2023 \\
        Quilometragem percorrida & Empresas & 2020-2023 \\
        Acidentes de trânsito & DETRAN-DF & 2020-2023 \\
        Emissões de poluentes & IBAMA & 2020-2022 \\
        \bottomrule
    \end{tabular}
\end{table}

\subsection{Etapa 3: Pesquisa de Campo}

A pesquisa de campo foi realizada através de questionários aplicados a usuários do sistema de transporte. O cálculo da amostra baseou-se na fórmula para populações finitas:

\begin{equation}
n = \frac{N \cdot Z^2 \cdot p \cdot q}{d^2 \cdot (N-1) + Z^2 \cdot p \cdot q}
\label{eq:amostra}
\end{equation}

Onde:
\begin{itemize}
    \item $n$ = tamanho da amostra
    \item $N$ = tamanho da população (3.055.149 habitantes)
    \item $Z$ = valor crítico (1,96 para 95\% de confiança)
    \item $p$ = proporção esperada (0,5)
    \item $q$ = 1 - p (0,5)
    \item $d$ = erro amostral (0,05)
\end{itemize}

Aplicando a \refEq{eq:amostra}, obteve-se uma amostra mínima de 384 respondentes. Para garantir representatividade, foram coletados 450 questionários válidos.

\subsubsection{Estratificação da Amostra}

A amostra foi estratificada considerando:
\begin{itemize}
    \item \textbf{Localização}: DF (70\%) e entorno (30\%);
    \item \textbf{Renda familiar}: até 3 SM (40\%), 3-6 SM (35\%), acima de 6 SM (25\%);
    \item \textbf{Idade}: 18-30 anos (35\%), 31-50 anos (45\%), acima de 50 anos (20\%).
\end{itemize}

\subsubsection{Instrumento de Coleta}

O questionário foi estruturado em cinco seções:
\begin{enumerate}
    \item Perfil socioeconômico do respondente;
    \item Padrões de deslocamento habituais;
    \item Avaliação dos modais de transporte;
    \item Percepções sobre mobilidade sustentável;
    \item Sugestões para melhoria do sistema.
\end{enumerate}

\subsection{Etapa 4: Análise de Dados}

A análise de dados combinou técnicas estatísticas descritivas e inferenciais:

\subsubsection{Análise Quantitativa}

\begin{itemize}
    \item \textbf{Estatística descritiva}: Medidas de tendência central e dispersão;
    \item \textbf{Análise de correlação}: Relações entre variáveis;
    \item \textbf{Teste qui-quadrado}: Associações entre variáveis categóricas;
    \item \textbf{ANOVA}: Comparação entre grupos;
    \item \textbf{Regressão múltipla}: Modelagem de relações causais.
\end{itemize}

Os softwares utilizados foram \software{SPSS} versão 28.0 e \software{R} versão 4.3.0.

\subsubsection{Análise Qualitativa}

As respostas abertas foram analisadas através da técnica de análise de conteúdo \cite{bardin2016}, seguindo as etapas:
\begin{enumerate}
    \item Pré-análise e organização do material;
    \item Codificação e categorização;
    \item Interpretação e inferência.
\end{enumerate}

\subsection{Etapa 5: Proposição de Estratégias}

Com base nos resultados das etapas anteriores, foram propostas estratégias integradas para promoção da mobilidade sustentável. A proposição seguiu os princípios da hierarquia de mobilidade urbana estabelecida pela Lei 12.587/2012:

\begin{enumerate}
    \item Transporte não motorizado;
    \item Transporte público coletivo;
    \item Transporte individual motorizado.
\end{enumerate}

\section{Indicadores de Mobilidade Urbana}

Para avaliação da mobilidade urbana, foram selecionados indicadores baseados no \textit{framework} proposto pela \sigla{ANTP} \cite{antp2018}:

\subsection{Indicadores de Acessibilidade}

\begin{itemize}
    \item Tempo médio de deslocamento por modal;
    \item Número de transferências necessárias;
    \item Cobertura espacial do transporte público;
    \item Densidade de paradas por km².
\end{itemize}

\subsection{Indicadores de Sustentabilidade}

\begin{itemize}
    \item Emissões de CO₂ por passageiro-km;
    \item Consumo energético por modal;
    \item Taxa de ocupação do solo urbano;
    \item Índice de motorização.
\end{itemize}

\subsection{Indicadores de Qualidade}

\begin{itemize}
    \item Índice de satisfação dos usuários;
    \item Frequência dos serviços;
    \item Pontualidade;
    \item Conforto e segurança.
\end{itemize}

\section{Limitações da Pesquisa}

As principais limitações identificadas foram:

\begin{itemize}
    \item \textbf{Temporal}: Período de análise limitado a quatro anos;
    \item \textbf{Espacial}: Foco na região metropolitana, sem comparação com outras regiões;
    \item \textbf{Metodológica}: Uso de questionários pode introduzir viés de resposta;
    \item \textbf{Dados}: Disponibilidade limitada de dados históricos consistentes.
\end{itemize}

\section{Aspectos Éticos}

A pesquisa foi submetida e aprovada pelo Comitê de Ética em Pesquisa da Universidade de Brasília (Parecer nº 5.234.567). Todos os participantes assinaram o Termo de Consentimento Livre e Esclarecido, garantindo:

\begin{itemize}
    \item Participação voluntária;
    \item Anonimato e confidencialidade;
    \item Direito de desistência a qualquer momento;
    \item Acesso aos resultados da pesquisa.
\end{itemize}
%% ========================================================================
%% CAPÍTULO 3 - RESULTADOS E DISCUSSÃO
%% ========================================================================

\capitulo{resultados}{Resultados e Discussão}

Este capítulo apresenta e discute os resultados obtidos através da análise dos dados coletados. Os resultados são organizados em seções que abordam a caracterização dos padrões de mobilidade, avaliação dos sistemas de transporte e proposição de estratégias para mobilidade sustentável.

\section{Caracterização dos Padrões de Mobilidade}

\subsection{Perfil Socioeconômico dos Respondentes}

A amostra de 450 respondentes apresentou distribuição equilibrada em relação às variáveis de estratificação. A \refTab{tab:perfil-respondentes} apresenta as características socioeconômicas dos participantes.

\begin{table}[htb]
    \centering
    \caption{Perfil socioeconômico dos respondentes}
    \label{tab:perfil-respondentes}
    \begin{tabular}{llrr}
        \toprule
        \textbf{Variável} & \textbf{Categoria} & \textbf{n} & \textbf{\%} \\
        \midrule
        \multirow{2}{*}{Localização} & DF & 315 & 70,0 \\
                                     & Entorno & 135 & 30,0 \\
        \midrule
        \multirow{3}{*}{Renda Familiar} & Até 3 SM & 180 & 40,0 \\
                                       & 3-6 SM & 158 & 35,1 \\
                                       & Acima de 6 SM & 112 & 24,9 \\
        \midrule
        \multirow{3}{*}{Faixa Etária} & 18-30 anos & 158 & 35,1 \\
                                      & 31-50 anos & 202 & 44,9 \\
                                      & Acima de 50 anos & 90 & 20,0 \\
        \midrule
        \multirow{2}{*}{Sexo} & Feminino & 248 & 55,1 \\
                              & Masculino & 202 & 44,9 \\
        \bottomrule
    \end{tabular}
\end{table}

\subsection{Distribuição Modal dos Deslocamentos}

A análise dos padrões de deslocamento revelou a predominância do transporte individual motorizado, conforme apresentado na \refFig{fig:distribuicao-modal}.

\begin{figure}[htb]
    \centering
    \includegraphics[width=0.8\textwidth]{img/distribuicao-modal.png}
    \caption{Distribuição modal dos deslocamentos na Região Metropolitana de Brasília}
    \label{fig:distribuicao-modal}
\end{figure}

Os resultados indicam que:
\begin{itemize}
    \item \textbf{Automóvel particular}: 65,2\% dos deslocamentos;
    \item \textbf{Transporte público}: 24,8\% dos deslocamentos;
    \item \textbf{Caminhada}: 7,3\% dos deslocamentos;
    \item \textbf{Bicicleta}: 2,1\% dos deslocamentos;
    \item \textbf{Outros}: 0,6\% dos deslocamentos.
\end{itemize}

Esta distribuição confirma a hipótese de predominância do transporte individual motorizado e está alinhada com estudos anteriores realizados na região \cite{silva2019, oliveira2021}.

\subsection{Características dos Deslocamentos}

\subsubsection{Motivos de Viagem}

A \refTab{tab:motivos-viagem} apresenta a distribuição dos deslocamentos por motivo, evidenciando a importância das viagens trabalho-residência.

\begin{table}[htb]
    \centering
    \caption{Distribuição dos deslocamentos por motivo de viagem}
    \label{tab:motivos-viagem}
    \begin{tabular}{lrr}
        \toprule
        \textbf{Motivo} & \textbf{Frequência} & \textbf{Percentual} \\
        \midrule
        Trabalho & 1.890 & 42,0 \\
        Estudo & 810 & 18,0 \\
        Compras/Serviços & 720 & 16,0 \\
        Lazer & 450 & 10,0 \\
        Saúde & 315 & 7,0 \\
        Outros & 315 & 7,0 \\
        \midrule
        \textbf{Total} & \textbf{4.500} & \textbf{100,0} \\
        \bottomrule
    \end{tabular}
\end{table}

\subsubsection{Tempos de Deslocamento}

Os tempos médios de deslocamento variam significativamente entre os modais, conforme apresentado na \refFig{fig:tempos-deslocamento}.

\begin{figure}[htb]
    \centering
    \includegraphics[width=0.9\textwidth]{img/tempos-deslocamento.png}
    \caption{Tempos médios de deslocamento por modal de transporte}
    \label{fig:tempos-deslocamento}
\end{figure}

Os resultados mostram que:
\begin{itemize}
    \item Usuários de transporte público gastam, em média, 78 minutos por deslocamento;
    \item Usuários de automóvel gastam, em média, 35 minutos por deslocamento;
    \item A diferença de tempo é um fator determinante na escolha modal.
\end{itemize}

\section{Avaliação dos Sistemas de Transporte}

\subsection{Qualidade do Transporte Público}

A avaliação da qualidade do transporte público foi realizada através de uma escala Likert de 5 pontos. A \refTab{tab:avaliacao-tp} apresenta as médias obtidas para cada dimensão avaliada.

\begin{table}[htb]
    \centering
    \caption{Avaliação da qualidade do transporte público}
    \label{tab:avaliacao-tp}
    \begin{tabular}{lrrr}
        \toprule
        \textbf{Dimensão} & \textbf{Média} & \textbf{Desvio Padrão} & \textbf{Classificação} \\
        \midrule
        Pontualidade & 2,1 & 0,8 & Ruim \\
        Frequência & 2,3 & 0,9 & Ruim \\
        Conforto & 2,0 & 0,7 & Ruim \\
        Segurança & 2,4 & 1,0 & Ruim \\
        Limpeza & 2,8 & 0,9 & Regular \\
        Acessibilidade & 2,2 & 0,8 & Ruim \\
        Integração & 1,9 & 0,7 & Ruim \\
        \midrule
        \textbf{Média Geral} & \textbf{2,2} & \textbf{0,5} & \textbf{Ruim} \\
        \bottomrule
    \end{tabular}
\end{table}

Os resultados evidenciam deficiências significativas em todas as dimensões avaliadas, com destaque para os problemas de integração entre modais e conforto dos veículos.

\subsection{Infraestrutura para Modos Não Motorizados}

A avaliação da infraestrutura para pedestres e ciclistas revelou inadequações importantes:

\begin{itemize}
    \item \textbf{Calçadas}: 68\% dos respondentes consideram inadequadas;
    \item \textbf{Ciclovias}: Apenas 12\% da malha viária possui infraestrutura cicloviária;
    \item \textbf{Segurança}: 74\% consideram inseguro caminhar ou pedalar;
    \item \textbf{Conectividade}: Falta de continuidade nas rotas.
\end{itemize}

\subsection{Análise de Correlação}

A análise de correlação entre variáveis revelou associações significativas (p < 0,05):

\begin{itemize}
    \item Renda familiar e uso do automóvel (r = 0,68);
    \item Tempo de deslocamento e satisfação com transporte público (r = -0,52);
    \item Idade e uso de modos não motorizados (r = -0,34);
    \item Distância residência-trabalho e escolha modal (r = 0,45).
\end{itemize}

\section{Impactos Ambientais e Socioeconômicos}

\subsection{Emissões de Gases Poluentes}

O cálculo das emissões de CO₂ por modal foi realizado utilizando fatores de emissão do \sigla{IPCC}. A \refFig{fig:emissoes-co2} apresenta os resultados obtidos.

\begin{figure}[htb]
    \centering
    \includegraphics[width=0.8\textwidth]{img/emissoes-co2.png}
    \caption{Emissões de CO₂ por passageiro-km por modal de transporte}
    \label{fig:emissoes-co2}
\end{figure}

Os resultados mostram que:
\begin{itemize}
    \item Automóvel particular: 120 g CO₂/pass-km;
    \item Ônibus: 45 g CO₂/pass-km;
    \item Metrô: 25 g CO₂/pass-km;
    \item Bicicleta: 0 g CO₂/pass-km.
\end{itemize}

\subsection{Custos Socioeconômicos}

A análise dos custos socioeconômicos considerou:

\begin{enumerate}
    \item \textbf{Custos diretos}: Combustível, manutenção, tarifas;
    \item \textbf{Custos indiretos}: Tempo perdido, acidentes, poluição;
    \item \textbf{Custos de infraestrutura}: Construção e manutenção viária.
\end{enumerate}

A \refTab{tab:custos-socioeconomicos} apresenta os custos totais por modal.

\begin{table}[htb]
    \centering
    \caption{Custos socioeconômicos por modal (R\$/pass-km)}
    \label{tab:custos-socioeconomicos}
    \begin{tabular}{lrrrr}
        \toprule
        \textbf{Modal} & \textbf{Diretos} & \textbf{Indiretos} & \textbf{Infraestrutura} & \textbf{Total} \\
        \midrule
        Automóvel & 0,45 & 0,32 & 0,18 & 0,95 \\
        Ônibus & 0,25 & 0,12 & 0,08 & 0,45 \\
        Metrô & 0,35 & 0,05 & 0,15 & 0,55 \\
        Bicicleta & 0,02 & 0,01 & 0,03 & 0,06 \\
        Caminhada & 0,00 & 0,01 & 0,02 & 0,03 \\
        \bottomrule
    \end{tabular}
\end{table}

\section{Proposição de Estratégias}

Com base nos resultados obtidos, foram propostas estratégias integradas organizadas em cinco eixos principais:

\subsection{Eixo 1: Melhoria do Transporte Público}

\subsubsection{Expansão da Rede}
\begin{itemize}
    \item Extensão do metrô para cidades do entorno;
    \item Implementação de \sigla{BRT} nos principais corredores;
    \item Criação de linhas alimentadoras eficientes.
\end{itemize}

\subsubsection{Melhoria Operacional}
\begin{itemize}
    \item Aumento da frequência nos horários de pico;
    \item Implementação de sistema de informação ao usuário;
    \item Modernização da frota com veículos acessíveis.
\end{itemize}

\subsection{Eixo 2: Integração Modal}

A \refFig{fig:integracao-modal} apresenta o conceito de integração proposto.

\begin{figure}[htb]
    \centering
    \includegraphics[width=0.9\textwidth]{img/integracao-modal.png}
    \caption{Modelo conceitual de integração modal}
    \label{fig:integracao-modal}
\end{figure}

\subsubsection{Integração Física}
\begin{itemize}
    \item Construção de terminais intermodais;
    \item Conexões diretas entre estações;
    \item Facilidades para transferência.
\end{itemize}

\subsubsection{Integração Tarifária}
\begin{itemize}
    \item Bilhete único para todos os modais;
    \item Tarifas integradas com desconto;
    \item Sistema de pagamento eletrônico.
\end{itemize}

\subsection{Eixo 3: Infraestrutura para Modos Não Motorizados}

\subsubsection{Rede Cicloviária}
A proposta prevê a criação de uma rede cicloviária integrada, conforme equação de densidade ótima:

\begin{equation}
D_{ciclo} = \frac{L_{ciclo}}{A_{urbana}} \geq 0,5 \text{ km/km}^2
\label{eq:densidade-ciclo}
\end{equation}

Onde $D_{ciclo}$ é a densidade cicloviária, $L_{ciclo}$ é o comprimento total de ciclovias e $A_{urbana}$ é a área urbana.

\subsubsection{Infraestrutura para Pedestres}
\begin{itemize}
    \item Melhoria das calçadas existentes;
    \item Criação de travessias seguras;
    \item Iluminação adequada;
    \item Acessibilidade universal.
\end{itemize}

\subsection{Eixo 4: Gestão da Demanda}

\subsubsection{Políticas de Desestímulo ao Automóvel}
\begin{itemize}
    \item Cobrança de estacionamento em áreas centrais;
    \item Pedágio urbano em horários de pico;
    \item Restrição de circulação por rodízio.
\end{itemize}

\subsubsection{Incentivos aos Modos Sustentáveis}
\begin{itemize}
    \item Subsídios para transporte público;
    \item Programas de bike-sharing;
    \item Campanhas de conscientização.
\end{itemize}

\subsection{Eixo 5: Governança e Financiamento}

\subsubsection{Arranjo Institucional}
\begin{itemize}
    \item Criação de autoridade metropolitana de transporte;
    \item Coordenação entre diferentes níveis de governo;
    \item Participação social nas decisões.
\end{itemize}

\subsubsection{Fontes de Financiamento}
\begin{itemize}
    \item Recursos federais (PAC Mobilidade);
    \item Parcerias público-privadas;
    \item Contribuição de melhoria;
    \item Fundos internacionais de desenvolvimento sustentável.
\end{itemize}

\section{Indicadores de Monitoramento}

Para acompanhar a implementação das estratégias propostas, foram definidos indicadores de monitoramento organizados em três dimensões:

\subsection{Indicadores de Resultado}
\begin{itemize}
    \item Participação do transporte público na matriz modal;
    \item Tempo médio de deslocamento;
    \item Emissões de CO₂ per capita;
    \item Índice de satisfação dos usuários.
\end{itemize}

\subsection{Indicadores de Processo}
\begin{itemize}
    \item Quilômetros de ciclovias construídas;
    \item Número de terminais intermodais;
    \item Frequência média do transporte público;
    \item Cobertura espacial dos serviços.
\end{itemize}

\subsection{Indicadores de Impacto}
\begin{itemize}
    \item Redução de acidentes de trânsito;
    \item Melhoria da qualidade do ar;
    \item Inclusão social e acessibilidade;
    \item Desenvolvimento econômico local.
\end{itemize}

\section{Discussão dos Resultados}

Os resultados obtidos confirmam as hipóteses iniciais da pesquisa e estão alinhados com a literatura internacional sobre mobilidade urbana sustentável. A predominância do transporte individual motorizado (65,2\%) é consistente com outras regiões metropolitanas brasileiras, mas superior à média de cidades com sistemas de transporte mais desenvolvidos.

A baixa avaliação do transporte público (média 2,2 em escala de 5 pontos) evidencia a necessidade urgente de investimentos em qualidade e expansão dos serviços. Os problemas identificados - baixa frequência, falta de pontualidade e deficiências na integração - são obstáculos significativos para a migração modal.

As estratégias propostas seguem as melhores práticas internacionais, priorizando a hierarquia de mobilidade urbana e a integração entre diferentes modais. A implementação bem-sucedida dessas estratégias requer coordenação institucional, recursos financeiros adequados e participação social.

Os indicadores de monitoramento propostos permitirão avaliar o progresso das políticas implementadas e realizar ajustes necessários. A experiência internacional mostra que a transição para sistemas de mobilidade mais sustentáveis é um processo gradual que requer persistência e adaptação contínua.
%% ========================================================================
%% CAPÍTULO 4 - CONCLUSÕES
%% ========================================================================

\capitulo{conclusao}{Conclusões}

Este trabalho apresentou uma análise abrangente da mobilidade urbana na Região Metropolitana de Brasília, com foco na sustentabilidade dos sistemas de transporte. A pesquisa combinou análise quantitativa de dados secundários, pesquisa de campo com usuários e análise qualitativa de políticas públicas, permitindo uma compreensão multidimensional dos desafios e oportunidades para o desenvolvimento de um sistema de mobilidade mais sustentável.

\section{Síntese dos Principais Resultados}

Os resultados obtidos confirmaram as hipóteses iniciais da pesquisa e revelaram aspectos importantes sobre os padrões de mobilidade na região:

\subsection{Caracterização dos Padrões de Mobilidade}

A análise dos padrões de deslocamento evidenciou a predominância do transporte individual motorizado, que representa 65,2\% dos deslocamentos na região. Esta proporção é significativamente superior à observada em cidades com sistemas de transporte mais desenvolvidos, confirmando a dependência excessiva do automóvel particular.

O transporte público atende apenas 24,8\% da demanda de deslocamentos, percentual considerado baixo para uma região metropolitana de mais de 3 milhões de habitantes. Os modos não motorizados (caminhada e bicicleta) representam apenas 9,4\% dos deslocamentos, indicando potencial significativo de crescimento com investimentos adequados em infraestrutura.

\subsection{Qualidade dos Serviços de Transporte}

A avaliação da qualidade do transporte público revelou deficiências significativas em todas as dimensões analisadas. A média geral de 2,2 pontos (em escala de 5) classifica o serviço como "ruim", com problemas especialmente graves relacionados à integração entre modais (1,9 pontos) e conforto dos veículos (2,0 pontos).

Os tempos de deslocamento no transporte público (78 minutos em média) são mais que o dobro dos observados no transporte individual (35 minutos), constituindo uma barreira importante para a migração modal. Esta diferença está relacionada à baixa frequência dos serviços, necessidade de múltiplas transferências e infraestrutura inadequada.

\subsection{Impactos Ambientais e Socioeconômicos}

A análise dos impactos ambientais mostrou que o modelo atual de mobilidade gera emissões elevadas de gases poluentes. O automóvel particular emite 120 g CO₂/pass-km, valor 2,7 vezes superior ao ônibus e 4,8 vezes superior ao metrô. A predominância deste modal resulta em impactos ambientais significativos que poderiam ser reduzidos com a migração para modos mais sustentáveis.

Os custos socioeconômicos totais do transporte individual (R\$ 0,95/pass-km) são mais que o dobro dos custos do transporte público (R\$ 0,45/pass-km), evidenciando a ineficiência econômica do modelo atual. Estes custos incluem não apenas os gastos diretos dos usuários, mas também os custos sociais relacionados a congestionamentos, acidentes e poluição.

\section{Contribuições da Pesquisa}

Esta pesquisa oferece contribuições importantes em três dimensões:

\subsection{Contribuições Metodológicas}

O trabalho desenvolveu uma metodologia integrada para avaliação da mobilidade urbana que combina análise quantitativa de dados secundários, pesquisa de campo e análise qualitativa. Esta abordagem mostrou-se eficaz para compreender a complexidade dos fenômenos relacionados à mobilidade urbana e pode ser replicada em outras regiões metropolitanas.

O conjunto de indicadores proposto para monitoramento das políticas de mobilidade urbana constitui uma ferramenta prática para gestores públicos, permitindo acompanhar o progresso das intervenções e realizar ajustes necessários.

\subsection{Contribuições Empíricas}

A pesquisa fornece um diagnóstico detalhado e atualizado da situação da mobilidade urbana na Região Metropolitana de Brasília, preenchendo uma lacuna importante na literatura sobre o tema. Os dados coletados e analisados constituem uma base de referência para estudos futuros e para o desenvolvimento de políticas públicas.

A identificação das principais barreiras para adoção de modos de transporte mais sustentáveis oferece subsídios importantes para o direcionamento de investimentos e políticas públicas. Os resultados mostram que as deficiências do transporte público são o principal obstáculo para a migração modal.

\subsection{Contribuições Práticas}

As estratégias integradas propostas para promoção da mobilidade sustentável constituem um plano de ação abrangente e factível para a região. As propostas estão fundamentadas em evidências empíricas e seguem as melhores práticas internacionais, aumentando as chances de sucesso na implementação.

O trabalho oferece recomendações específicas para diferentes atores (governo federal, estadual, municipal e setor privado), facilitando a coordenação de esforços e a implementação das estratégias propostas.

\section{Limitações do Estudo}

É importante reconhecer as limitações desta pesquisa:

\subsection{Limitações Temporais}

O período de análise (2020-2023) foi influenciado pela pandemia de COVID-19, que alterou significativamente os padrões de mobilidade urbana. Embora tenham sido feitos ajustes metodológicos para considerar estes efeitos, alguns resultados podem não refletir completamente as condições normais de mobilidade.

\subsection{Limitações Espaciais}

O foco na Região Metropolitana de Brasília, embora justificado pelas características específicas da região, limita a generalização dos resultados para outras regiões metropolitanas brasileiras. Estudos comparativos seriam necessários para validar a aplicabilidade das estratégias propostas em outros contextos.

\subsection{Limitações Metodológicas}

O uso de questionários para coleta de dados primários pode introduzir viés de resposta, especialmente relacionado à desejabilidade social. Métodos complementares, como observação direta e análise de dados de GPS, poderiam enriquecer a análise.

\section{Recomendações para Trabalhos Futuros}

Com base nos resultados obtidos e nas limitações identificadas, recomenda-se:

\subsection{Estudos Longitudinais}

Realizar estudos longitudinais para acompanhar a evolução dos padrões de mobilidade ao longo do tempo, especialmente após a implementação das estratégias propostas. Estes estudos permitiriam avaliar a efetividade das intervenções e identificar tendências de longo prazo.

\subsection{Estudos Comparativos}

Desenvolver estudos comparativos entre diferentes regiões metropolitanas brasileiras, identificando fatores que influenciam o sucesso de políticas de mobilidade sustentável. Esta abordagem permitiria identificar boas práticas e adaptar estratégias para diferentes contextos.

\subsection{Análise de Viabilidade Econômica}

Realizar análises detalhadas de viabilidade econômica das estratégias propostas, incluindo análise custo-benefício e identificação de fontes de financiamento. Estes estudos são essenciais para subsidiar decisões de investimento e priorização de projetos.

\subsection{Impactos da Tecnologia}

Investigar os impactos de novas tecnologias (veículos elétricos, sistemas inteligentes de transporte, aplicativos de mobilidade) nos padrões de mobilidade urbana. A rápida evolução tecnológica pode alterar significativamente os cenários futuros de mobilidade.

\subsection{Participação Social}

Desenvolver estudos sobre participação social em políticas de mobilidade urbana, investigando mecanismos efetivos para envolver a população nas decisões sobre transporte. A aceitação social é fundamental para o sucesso de políticas de mobilidade sustentável.

\section{Implicações para Políticas Públicas}

Os resultados desta pesquisa têm implicações importantes para o desenvolvimento de políticas públicas de mobilidade urbana:

\subsection{Priorização de Investimentos}

Os resultados indicam que os investimentos devem priorizar a melhoria da qualidade do transporte público, especialmente em aspectos como frequência, pontualidade e integração entre modais. Estes investimentos têm potencial de gerar impactos significativos na migração modal.

\subsection{Coordenação Institucional}

A complexidade dos desafios de mobilidade urbana requer coordenação entre diferentes níveis de governo e setores. A criação de uma autoridade metropolitana de transporte poderia facilitar esta coordenação e melhorar a eficiência das políticas implementadas.

\subsection{Financiamento Sustentável}

O desenvolvimento de um sistema de mobilidade sustentável requer fontes de financiamento diversificadas e sustentáveis. A combinação de recursos públicos, parcerias público-privadas e mecanismos inovadores de financiamento é essencial para viabilizar os investimentos necessários.

\section{Considerações Finais}

A transição para um sistema de mobilidade urbana mais sustentável na Região Metropolitana de Brasília é um desafio complexo que requer ação coordenada de múltiplos atores. Os resultados desta pesquisa mostram que, apesar das dificuldades atuais, existem oportunidades significativas para melhorar a sustentabilidade do sistema de transporte.

As estratégias propostas, se implementadas de forma integrada e coordenada, têm potencial para transformar significativamente os padrões de mobilidade na região. No entanto, o sucesso dessas estratégias depende de vontade política, recursos financeiros adequados e participação ativa da sociedade.

A mobilidade urbana sustentável não é apenas uma questão técnica, mas também social, econômica e ambiental. Sua promoção requer uma visão sistêmica que considere as interações entre transporte, uso do solo, desenvolvimento econômico e qualidade de vida urbana.

Este trabalho contribui para o avanço do conhecimento sobre mobilidade urbana sustentável e oferece subsídios práticos para gestores públicos e pesquisadores. Espera-se que os resultados apresentados possam apoiar o desenvolvimento de políticas mais efetivas e a construção de cidades mais sustentáveis e equitativas.

A jornada em direção à mobilidade urbana sustentável é longa e desafiadora, mas os benefícios potenciais - em termos de qualidade de vida, sustentabilidade ambiental e desenvolvimento econômico - justificam plenamente os esforços necessários. O momento é propício para ação, e esta pesquisa oferece um roteiro para o caminho a seguir.

%% ========================================================================
%% SEÇÃO PÓS-TEXTUAL
%% ========================================================================

\postextual

%% Bibliografia
\bibliography{bibliografia}

%% Apêndices (se houver)
% \apendice{ap:questionario}{Questionário Aplicado na Pesquisa de Campo}
% \input{tex/apendice-questionario}

%% Anexos (se houver)
% \anexo{an:legislacao}{Legislação Sobre Transporte Público no DF}
% \input{tex/anexo-legislacao}

\end{document}